In this chapter we describe how language models deal with multimodality.
By multimodality we mean information which is outside of the main input text.
Prototypically, the extra modality is visual.
For this, consider the case of an image captioning model.
It is a language model which also has a module that is able to process an image and generate text based on it.
However, there are also other types of multimodalities, which we cover in the subsequent section.
These include, for example, knowledge-graphs or knowledge-bases in general and we discuss how they can be used to enhance the model with extra knowledge.